\par\noindent\textbf{(1)} 从 \(a_i=-\partial_i\Phi=-GMx_i/r^3\) 计算空间导数，得到相邻质点 \(\delta a_i=(\partial_ja_i)\xi_j=\mathsf T_{ij}\xi_j\)，其中 \(\mathsf T_{ij}=GM(3x_ix_j/r^5-\delta_{ij}/r^3)\)。这一矩阵是引力加速度的真实梯度，而非某个旋转系中的离心项。
\par\noindent\textbf{(2)} 记径向单位矢量 \(n_i=x_i/r\)，则 \(\mathsf T=(GM/r^3)(3\boldsymbol n\boldsymbol n^T-I)\)。因此 \(\mathsf T\boldsymbol n=(2GM/r^3)\boldsymbol n\)，任意 \(\boldsymbol e\perp\boldsymbol n\) 有 \(\mathsf T\boldsymbol e=-(GM/r^3)\boldsymbol e\)。径向拉伸而横向压缩，三个本征值为 \(2GM/r^3,-GM/r^3,-GM/r^3\)。
\par\noindent\textbf{(3)} 本征值之和 \(\operatorname{tr}\mathsf T=0\) 与真空 \(\nabla^2\Phi=0\) 一致。对均匀密度实心球内部，\(\nabla^2\Phi=4\pi G\rho\)，故 \(\operatorname{tr}\mathsf T=-4\pi G\rho\ne0\)，点质量真空潮汐表达式不再适用。边界应通过连续引力势与场匹配。
\par\noindent\textbf{(4)} 半径 \(a\)、密度 \(\rho_s\) 的小天体表面自引力约为 \(g_s=4\pi G\rho_s a/3\)；其两端潮汐差加速度量级 \(\Delta g\simeq2GM a/r^3\)。令两者同阶得 \(r_R\sim[3M/(2\pi\rho_s)]^{1/3}\)。该尺度只是一阶估计，精确Roche极限还取决于潮汐响应、同步转动、流体形状与材料强度。
